Este livro nasceu a partir das transcrições das aulas da disciplina de Programação Web. Diferentemente de outros materiais de consulta mais enxutos, este livro foi pensado para ser um pouco mais completo: além de organizar e explicar o conteúdo apresentado em aula, ele traz exemplos de código em HTML, CSS, JavaScript e Java/Spring, além de diagramas explicativos para os conceitos mais visuais, como o modelo de caixa do CSS, a árvore do DOM e a arquitetura cliente-servidor. A estrutura do livro segue a mesma sequência lógica adotada durante as aulas: parte-se dos fundamentos da profissão e da web (redes, protocolos, ferramentas, hospedagem), avança-se para o front-end com HTML e CSS, passa-se pelos fundamentos de JavaScript e, por fim, apresenta-se o back-end com Java, Spring Boot, Docker e persistência de dados. Como esta é uma disciplina que o leitor ainda cursará, com foco especial na parte de front-end — área na qual ele parte do zero —, alguns capítulos foram enriquecidos com conteúdo complementar não apresentado originalmente nas aulas (por exemplo, estruturas de controle e funções em JavaScript, e a manipulação do DOM), de modo a preencher lacunas naturais de um curso introdutório e tornar a leitura mais autossuficiente. Esses complementos estão sempre claramente integrados ao restante do texto, mas a base principal e a linha condutora do conteúdo continuam sendo o material e a fala do professor responsável pela disciplina.

Boa leitura e bons estudos.
